\DTMsavedate{mydate}{2022-06-30}
\maketitle{Industrial Control}{Electronic Engineering}{\DTMusedate{mydate}}

% ----------------------------------------------------- %
% COURSE INFO
\subsection*{Course Information}\label{course:5241}
\vspace{0ex}%
\noindent\parbox{0.5\textwidth}{%
    \noindent\begin{tabular}{@{}ll}
                 \textsf{Course:}          & 	Industrial Control  \\
                 \textsf{Course Code:}         & 	TEE 5241    \\
                 \textsf{Credit Hours:}    & 	10
    \end{tabular}}
\vspace{2ex}\hrule\vspace{2ex}

% ----------------------------------------------------- %
% INSTRUCTOR INFO
\subsection*{Lecturer Information}
\noindent\parbox{0.5\textwidth}{%
    \noindent\begin{tabular}{@{}ll}
                 \textsf{Name:}         &     Reginald Gonye         \\
%\textsf{Phone:}        &    \phone          \\
\textsf{Email:}        &    reginald.gonye@nust.ac.zw \\
\textsf{Qualifications:}        &  MPhil in Electronic Engineering, (Hon) BEng in Electronic Engineering
    \end{tabular}}
\vspace{2ex}\hrule\vspace{2ex}


% ----------------------------------------------------- %
% COURSE SYNOPSIS
\subsection*{Course Synopsis}
The module focuses on industrial control situations, process control, instrumentation, actuators, transducers and controllers, hybrid systems, time-domain analysis, state-space analysis, stability; computer control, system characterization, algorithm design, feedback control for digital systems and PLC applications.
\vspace{2ex} \hrule\vspace{2ex}


% ----------------------------------------------------- %
% COURSE DESCRIPTION
\subsection*{Learning Aim}
To introduce students to practical arrangements of Control Systems in industry.
To instill skills of developing PLC based solutions for control situations into students
\vspace{2ex} \hrule\vspace{2ex}

% ----------------------------------------------------- %
% LEARNING OBJECTIVES
\subsection*{Learning Objectives}
By the end of the course, candidates should be able to:
\begin{outline}
    \1      Understand the arrangement of typical control systems in industry.

    \1  	Understand the operation of a PLC.

    \1  	Be able to develop PLC-based solutions for various control situations.

\end{outline}
\vspace{2ex}\hrule\vspace{2ex}

% ----------------------------------------------------- %
% COURSE TOPICS
\subsection*{Topics}
\begin{outline}
    \1 A review of Control Systems
    \1 Sensors and Transducers
    \1 Actuators and their Control
    \1 Computers in Industrial Control
    \1 The PLC
    \1 SCADA
\end{outline}
\vspace{2ex} \hrule\vspace{2ex}

% ----------------------------------------------------- %
% Students Assenssment
\subsection*{Assessment of Students}
\begin{outline}
    \1 Continuous assessment will consist of tests and assignments (NB: Attendance may also contribute to the coursework mark).
    \1 Final assessment will consist of an examination at the end of the semester. The examination shall comprise of six questions each carrying 25 marks. You will be required to attempt four of the six questions.
\end{outline}
\vspace{2ex}\hrule\vspace{2ex}

% ----------------------------------------------------- %
% Grade Evaluation
\subsection*{Course Evaluation and Grading Policies}
Grades are based on:
Continuous assessment (\qty{25}{\percent}),
Final Exam (\qty{75}{\percent})
\vspace{2ex}\hrule
\vspace{2ex}\hrule\vspace{2ex}
% ----------------------------------------------------- %
% EQUIPMENT
% https://sites.udel.edu/computing-purchases/personal-specs/
% \subsection*{Essential Equipment}

% \vspace{2ex}\hrule\vspace{2ex}

% Policies and Other information
%\subsection*{Policies and Other Information}
% Conduct
%\subsection*{Key Text}

%\begin{itemize}
%    \item Rafiquzzaman Mohamed, Microprocessor Theory and Applications with 68000/68020 and Pentium, 2008 (John Wiley & Sones, INc, Hobolen, New Jersey/Canada)
%\end{itemize}
